% CHAPTER 4: METHODOLOGY
% Auto-generated from codebase inspection of the deepfake-project repository.
\documentclass[11pt]{report}
\usepackage[utf8]{inputenc}
\usepackage{graphicx}
\usepackage{booktabs}
% (file replaced with expanded academic-style content preserving implementation facts)
\documentclass[11pt]{report}
\usepackage[utf8]{inputenc}
\usepackage{graphicx}
\usepackage{booktabs}
\usepackage{hyperref}
\usepackage{amsmath}
\usepackage{longtable}
\begin{document}
\chapter*{CHAPTER 4: METHODOLOGY}
\addcontentsline{toc}{chapter}{CHAPTER 4: METHODOLOGY}

\section*{4.0 Methodology Overview}
This chapter details the methodology employed to construct, train and evaluate
the deepfake detection system implemented in this repository. The system
integrates multiple complementary signal domains: spatial RGB appearance,
handcrafted high-frequency residuals (SRM), and frequency-domain magnitude
features (FFT). These modalities are fused with modern deep feature extractors
and lightweight attention/pooling modules. The pipeline comprises: dataset
construction; preprocessing and augmentation; feature extraction using
pretrained backbones; optional preprocessing modules (SRM and FFT); attention
and pooling; optimisation and training; and evaluation. Implementation details
are referenced directly to repository modules (for example
	exttt{scripts/preprocessing/srm.py}, \texttt{scripts/training/train_full.py}).

The overarching design objective is to produce a classifier that reliably
discriminates among three semantic classes: Real, AI Generated, and AI Edited.
To this end, the methodology emphasises (a) careful dataset curation to avoid
leakage and ensure representativeness, (b) multi-domain preprocessing that
amplifies forensic cues, (c) transfer learning from ImageNet-pretrained
backbones to leverage robust low-level features, and (d) attention and pooling
mechanisms that concentrate model capacity on informative spatial and channel
components. The following sections expand each part of the pipeline while
preserving all factual implementation details present in the codebase.

\section*{4.1 Dataset Description}
4.1.1 Motivation for a Curated Dataset
\newline
Supervised forensic models are sensitive to dataset composition. Publicly
available collections may be biased towards a narrow set of generative models
or capture conditions; therefore the project implements a purpose-built,
multi-source dataset assembled by the pipeline under \texttt{dataset_builder/}.
This curation enables deterministic pre-processing, enforced class-balancing,
and provenance metadata that supports reproducible experiments and rigorous
error analysis.

4.1.2 Three-class Taxonomy
\newline
The dataset uses an explicit three-class taxonomy:
\begin{itemize}
  \item \textbf{Real} (class 0): original, unmanipulated images.
  \item \textbf{AI Generated} (class 1): fully synthetic images produced by
  generative models (GANs, diffusion models, etc.).
  \item \textbf{AI Edited} (class 2): real images that have been edited (for
  example via inpainting, face-swap, or localized retouching).
\end{itemize}

The separation of generated and edited classes enables the classifier to learn
distinct signatures: global spectral fingerprints often indicate fully
generated imagery, while localized residual anomalies better characterise
editing operations.

4.1.3 Dataset Size, Splits and Export
\newline
Recorded counts (see \texttt{DATASET.md}) are: real = 26,000; ai_generated =
26,000; ai_edited = 25,865; total = 77,865. The pipeline writes exported
splits to \texttt{dataset_builder/train/}, \texttt{dataset_builder/val/} and
	exttt{dataset_builder/test/}. While nominal YAML-configured ratios exist
(70/15/15), the implemented cluster-aware splitter produces an empirical
export approximating 40/30/30 in some runs; the split implementation is in
	exttt{dataset_builder/modules/splitter.py} and is orchestrated by
	exttt{dataset_builder/pipeline.py}.

4.1.4 Class Balancing and Quota Sampling
\newline
The production pipeline enforces per-source quotas and overall class balance to
prevent dominance of any single class or source. The sampling logic is in
	exttt{dataset_builder/modules/sampler.py} and deterministic behaviour is
ensured via \texttt{set_global_seed} in \texttt{dataset_builder/pipeline.py}.
Balanced datasets reduce training bias and ensure that per-class metrics are
meaningful for operational evaluation.

4.1.5 Source Diversity and Manipulation Types
\newline
The dataset aggregates diverse capture conditions and manipulation methods to
improve generalisation. Diversity includes multiple camera models, compression
artifacts, and generative model variants; manipulation types include full
generation, inpainting, and face-swapping. The pipeline preserves provenance
metadata (for example \texttt{export_index.csv}) to support stratified
analyses by source or manipulation family.

4.1.6 Deduplication, Validation and Dataset Quality
\newline
Deduplication is performed with perceptual hashing and validation stages check
resolution, blur and metadata (see \texttt{dataset_builder/modules/deduplicator.py}
and \texttt{dataset_builder/modules/validator.py}). Cluster-aware splitting
ensures that related items (e.g. frames extracted from the same video) remain
in the same split to prevent leakage. These steps are critical because low-
quality or duplicated training examples encourage memorisation of source
artifacts rather than learning transferable forensic cues.

\section*{4.2 Data Preprocessing}
4.2.1 Overview
\newline
Preprocessing responsibilities are split across dataset-level filtering and
model-input preparation. Model-input preprocessing is implemented primarily in
	exttt{scripts/preprocessing/preprocessing.py} and uses Albumentations for
training-time augmentations and Torchvision transforms for deterministic
inference-time processing. Optional preprocessing modules (SRM, FFT) are
implemented in \texttt{scripts/preprocessing/srm.py} and
	exttt{scripts/preprocessing/fft.py} and are applied via the
	exttt{PreprocessNet} wrapper in \texttt{scripts/training/train_full.py}.

4.2.2 Image Loading and Channel Ordering
\newline
At data ingest, images are read with OpenCV (BGR format) within
	exttt{scripts/dataloader/dataset.py}. The loader performs an explicit BGR→RGB
conversion prior to augmentation to align with pretrained backbones' expected
channel ordering and avoid inadvertent color-space anomalies.

4.2.3 Resizing and Normalization
\newline
All inputs are resized to 224×224 during augmentation. Pixel normalisation uses
ImageNet mean and standard deviation (mean=[0.485,0.456,0.406],
std=[0.229,0.224,0.225]) to maintain consistency with pretrained feature
initialisations. Conversion to PyTorch channel ordering (C,H,W) and floating
tensor is performed by \texttt{ToTensorV2} as the final augmentation step.

4.2.4 Augmentation Modes and Techniques
\newline
Training augmentations are exposed as three named transforms via
	exttt{get_train_transform(mode)}:
\begin{itemize}
  \item \texttt{light}: Resize(224,224), HorizontalFlip(p=0.5), Normalize(),
  ToTensorV2().
  \item \texttt{standard} (default): Resize(224,224), HorizontalFlip(p=0.5),
  Rotate(limit=15,p=0.5), RandomBrightnessContrast(p=0.5),
  ImageCompression(quality_range=(30,100),p=0.4), Normalize(), ToTensorV2().
  \item \texttt{strong}: as \texttt{standard} with intensified ImageCompression
  (20–95), additive noise (GaussNoise or ISONoise) and light blur (Gaussian or
  Median) to simulate harsh capture conditions.
\end{itemize}

Rationale: Compression simulation addresses lossy storage/transmission; noise
injection models sensor or post-processing noise; geometric transforms improve
robustness to framing variations; blur reduces sensitivity to fine-grained
texture so models learn coarser but more generalisable cues.

4.2.5 Training vs Inference Preprocessing
\newline
Training uses stochastic Albumentations to promote invariance, while
inference applies deterministic Torchvision transforms (Resize, CenterCrop,
ToTensor, Normalize) to ensure repeatable inputs and stable evaluation.

\section*{4.3 Feature Extraction}
4.3.1 Candidate Backbones and Feature Dimensions
\newline
The codebase supports multiple backbone families (constructed in
	exttt{scripts/training/train_full.py}): ResNet-18/ResNet-50 (final features
512/2048), ConvNeXt-Tiny/Small (final features 768), EfficientNet-B3 and
ViT-B/16 (supported with special handling for token pooling). Using
ImageNet-pretrained weights provides robust low-level filters and accelerates
convergence when fine-tuning for forensic tasks.

4.3.2 Convolutional Inductive Biases for Forensics
\newline
Convolutional architectures are well-suited for forensic analysis due to their
local receptive fields and hierarchical composition. Early layers capture
edges and textures while deeper layers encode increasingly semantic
representations; this hierarchy supports detection of both local manipulations
(inpainting seams) and more global inconsistencies (colour mismatches or
artificial textures). Transformer-based backbones are supported to compare
their global attention properties, but their tokenisation implies different
trade-offs.

4.3.3 Pooling and Classification Head
\newline
After spatial aggregation (GeM or average pooling) the pooled vector is
subjected to dropout (default 0.4) and then passed to a linear classification
head mapping to three logits. The head construction is performed by
	exttt{make_head(in_features)} in \texttt{scripts/training/train_full.py} and
is tuned to each backbone's final feature dimensionality.

\section*{4.4 Preprocessing Feature Modules}
4.4.1 SRM (Steganalysis Rich Model)
\newline
Implementation: \texttt{SRMLayer} (\texttt{scripts/preprocessing/srm.py})
computes three fixed high-pass residual channels using handcrafted kernels
(Laplacian 5×5, horizontal 5×5, vertical 5×5) applied to a grayscale image;
residuals are clamped (default clamp range 2.0) and concatenated with RGB to
form a 6-channel input. The helper \texttt{adapt_conv1_for_srm} modifies the
backbone's first convolution to accept 6 channels by initialising additional
weights to 0.1× the pretrained RGB weights.

Justification: SRM residuals amplify small-scale discontinuities and blending
seams that are often obscured by colour and texture in the RGB domain. By
presenting these residuals explicitly to the network, early convolutions can
detect structural anomalies that are indicative of manipulation.

4.4.2 FFT (Frequency Magnitude)
\newline
Implementation: \texttt{FFTLayer} (\texttt{scripts/preprocessing/fft.py})
computes the log-magnitude of the centred 2D Fourier transform of the
grayscale image, normalises it to [0,1], and returns a B×1×H×W tensor. When
enabled, this channel is concatenated after RGB and SRM channels (rgb →
residuals → fft) and the backbone's first convolution is adapted accordingly.

Justification: Frequency-domain analysis reveals global periodicities and
power-spectrum anomalies often characteristic of generative models or
resampling introduced by editing tools. FFT magnitude provides complementary
evidence to SRM and RGB features and can distinguish manipulation classes
that are ambiguous in the spatial domain.

\section*{4.5 Attention and Pooling Mechanisms}
4.5.1 CBAM: Channel and Spatial Attention
\newline
CBAM is implemented in \texttt{scripts/modules/attention_heads.py} and applies
channel attention (via pooled statistics and an MLP) followed by spatial
attention (via concatenated avg/max maps and a convolution). Applied to late
feature maps, CBAM reweights channels and spatial locations to emphasise
forensically salient information.

4.5.2 GeM Pooling
\newline
GeM is implemented as \texttt{GeM} (\texttt{scripts/modules/attention_heads.py}).
Given activation map $X \in \mathbb{R}^{C\times H\times W}$ GeM computes:
$$
\mathrm{GeM}(X)_c = \left(\frac{1}{HW} \sum_{i=1}^{H} \sum_{j=1}^{W} X_{c,i,j}^{p} \right)^{1/p}.
$$
The parameter $p$ (default 3.0) is registered as a tensor and may be frozen
or learned. GeM interpolates between average ($p=1$) and max pooling
($p\to\infty$), giving stronger emphasis to local high activations which are
useful when small regions contain manipulation evidence.

4.5.3 Role of Attention and GeM in Forensics
\newline
Attention reduces the influence of irrelevant background structure and lets the
model focus on high-signal regions. GeM ensures that strong local cues are
preserved during aggregation rather than being diluted by large spatial
averages; this property is particularly beneficial when manipulations occupy a
small fraction of the image.

\section*{4.6 Model Architecture}
4.6.1 Input and PreprocessNet Wrapper
\newline
The input pipeline optionally applies \texttt{SRMLayer} and/or \texttt{FFTLayer}
to an RGB tensor and concatenates channels in the order rgb → residuals → fft.
The concatenated tensor is provided to an adapted backbone via the
	exttt{PreprocessNet} wrapper implemented in \texttt{scripts/training/train_full.py}.

4.6.2 Backbone, Attention and Pooling
\newline
The adapted backbone (ResNet/ConvNeXt/EfficientNet/ViT) computes hierarchical
features. CBAM blocks may be inserted after late stages (for ResNet, after
	exttt{layer4}; for ConvNeXt, after \texttt{features}) and GeM may replace the
global average pooling layer. The resulting pooled vector is passed through a
dropout layer and a linear classification head (softmax for probabilities).

4.6.3 Step-by-step Data Flow
\begin{enumerate}
  \item Image file loaded (cv2) and converted BGR→RGB.
  \item Training augmentations applied (Resize → geometric → photometric →
  Normalize → ToTensorV2) or deterministic inference transforms used.
  \item SRM / FFT modules compute additional channels if enabled and outputs
  are concatenated with RGB.
  \item The concatenated tensor is forwarded through the adapted backbone, with
  optional CBAM applied to late feature maps.
  \item GeM or average pooling aggregates spatial maps to a vector, followed by
  dropout and the final linear head producing three logits.
\end{enumerate}

\section*{4.7 Training Procedure}
4.7.1 Losses and Class-weighting
\newline
The training code implements CrossEntropy, class-weighted CrossEntropy,
Focal Loss (\texttt{FocalLoss}) and a factory \texttt{build_criterion} for
WeightedFocalLoss. Default class weights are \texttt{[1.5, 1.0, 1.5]}. The
default focal gamma in code is 3.0; however model cards and many experiments
use $\gamma=2.0$ in practice. Label smoothing (commonly 0.1) is supported to
reduce overconfidence.

Rationale: Focal loss focuses optimisation on hard examples, which is useful
when manipulated examples are subtle. Class weights compensate for residual
sampling imbalance and emphasize operationally critical classes.

4.7.2 Optimisers, Schedules and Hyperparameters
\newline
Optimisers: Adam (default) and SGD (option). Default weight decay: 1e-4. The
learning rate default is 1e-4. Schedulers: CosineAnnealingWarmRestarts
(\texttt{T_0=20}, \texttt{T_mult=2}) for \texttt{lr_schedule='cosine'}, or
ReduceLROnPlateau for \texttt{lr_schedule='plateau'}. Batch size default is 64
and default epochs is 50 (baseline scripts use smaller defaults for quick
iterations). Early stopping and checkpointing save the best model (for example
to \texttt{best_model.pth}) and per-epoch metrics are written to
	exttt{metrics.csv} and \texttt{training_summary.json}.

Rationale: Cosine restarts provide periodic learning rate resets that help
escape local minima during long runs; Adam provides adaptive gradient steps that
are robust to noisy gradients introduced by augmentations.

4.7.3 Mixed Precision and Performance
\newline
Training uses AMP (\texttt{torch.amp.GradScaler}, \texttt{autocast}) and
enables \texttt{torch.compile} when available. These pragmatics increase
throughput and reduce memory footprint, enabling larger batch sizes and
faster experimentation.

\section*{4.8 Evaluation Metrics}
4.8.1 Metrics Implemented
\begin{itemize}
  \item \textbf{Accuracy}: overall fraction correct; logged as
  	exttt{train_acc} and \texttt{val_acc}.
  \item \textbf{Precision / Recall}: computed per-class; precision measures the
  correctness of positive predictions while recall measures completeness of
  detection for the true positives.
  \item \textbf{F1 score}: harmonic mean of precision and recall; code computes
  macro and per-class F1 using \texttt{sklearn.metrics.f1_score}.
  \item \textbf{Confusion matrix}: reconstructed from saved logits and labels to
  provide detailed misclassification patterns.
\end{itemize}

4.8.2 ROC-AUC and Offline Analysis
\newline
ROC-AUC is not computed by default within the training scripts but can be
derived offline from saved logits and labels. ROC-AUC is threshold-independent
and especially useful when selecting operational thresholds that balance false
positives and false-negatives.

4.8.3 Complementary Metric Usage
\newline
Accuracy gives a single-number summary but can conceal class-specific failures;
precision, recall and F1 provide per-class diagnostic insight, while the
confusion matrix reveals asymmetric confusions that guide dataset curation and
loss reweighting.

\section*{4.9 Methodology Summary}
This chapter has presented a principled methodology integrating dataset
engineering, robust preprocessing, multi-domain feature extraction, attention
and pooling mechanisms, and carefully chosen optimisation strategies. The
approach is implemented in the repository with explicit modules to ensure
reproducibility. Each design decision is motivated by forensic considerations:
SRM emphasises high-frequency residuals for local edit detection; FFT reveals
global spectral fingerprints; GeM and CBAM concentrate model capacity on
salient, high-signal regions; and augmentation plus class-aware sampling
promote robust generalisation. The subsequent chapters evaluate this
methodology empirically and discuss limitations and avenues for future work.

\section*{DIAGRAMS}
The following diagrams are recommended for the thesis; each includes a title,
components to show, directional data-flow arrows, and an explanatory note.

\begin{enumerate}
  \item \textbf{Overall System Pipeline}
    \begin{itemize}
      \item Components: raw sources → indexing → validation → deduplication
      → sampling/balancing → cluster-aware splitter → exported train/val/test
      → training (PreprocessNet → Backbone → Head) → evaluation → model
      cards/metrics.
      \item Arrows: left-to-right, annotate intermediate artifacts such as
      	exttt{index.csv} and \texttt{export_index.csv}.
      \item Explanation: end-to-end reproducible workflow from raw data to
      published experiment artifacts.
    \end{itemize}

  \item \textbf{Dataset Preprocessing Pipeline}
    \begin{itemize}
      \item Components: indexing → validation checks (resolution/blur) →
      deduplication (pHash) → quota sampling → cluster-aware split → export.
      \item Arrows: show conditional branches (invalid → rejected) and
      deterministic seeding.
      \item Explanation: clarifies steps that ensure dataset quality and avoid
      leakage.
    \end{itemize}

  \item \textbf{CNN Backbone Architecture}
    \begin{itemize}
      \item Components: input → stem convs → residual/ConvNeXt blocks → stage
      outputs → pooling (Avg/GeM) → dropout → FC head. Include feature sizes
      (ResNet-18: 512; ConvNeXt: 768; ResNet-50: 2048).
      \item Arrows: show spatial downsampling and channel dimension changes.
      \item Explanation: hierarchical feature extraction for forensic signals.
    \end{itemize}

  \item \textbf{SRM Filter Integration}
    \begin{itemize}
      \item Components: RGB → grayscale → three SRM kernels (Laplacian,
      horizontal, vertical) → residual maps → concatenation with RGB →
      adapted conv1 (0.1× init for extra channels).
      \item Arrows: highlight concatenation and how SRM buffers are stored in
      checkpoints.
      \item Explanation: how residual channels are fused at model input.
    \end{itemize}

  \item \textbf{FFT Feature Extraction}
    \begin{itemize}
      \item Components: grayscale → FFT → fftshift → log-magnitude →
      normalise → FFT channel concatenated with RGB/SRM.
      \item Arrows: indicate transformation from spatial to frequency domain
      and back to a spatially-aligned magnitude image.
      \item Explanation: documents spectral fingerprint extraction.
    \end{itemize}

  \item \textbf{Attention Module (CBAM)}
    \begin{itemize}
      \item Components: input feature map → channel attention (avg/max pools →
      MLP → sigmoid) → multiplication → spatial attention (concat avg/max →
      conv → sigmoid) → multiplication → output.
      \item Arrows: sequential application channel→spatial.
      \item Explanation: how attention reweights channels and locations.
    \end{itemize}

  \item \textbf{Feature Extraction Pipeline}
    \begin{itemize}
      \item Components: RGB → (optional SRM / FFT) → adapted backbone →
      attention → pooling (GeM) → feature vector → classifier head.
      \item Arrows: show multi-domain fusion and final prediction flow.
      \item Explanation: end-to-end signal fusion for detection.
    \end{itemize}

  \item \textbf{Training Workflow}
    \begin{itemize}
      \item Components: DataLoader (Albumentations) → mixed-precision
      forward/backward → optimizer step → LR scheduler → checkpointing and
      logging to \texttt{metrics.csv} and TensorBoard.
      \item Arrows: iterate per-epoch with checkpoint outputs.
      \item Explanation: operational training loop and logging artifacts.
    \end{itemize}

  \item \textbf{Evaluation Pipeline}
    \begin{itemize}
      \item Components: test data loader → model inference (with SRM enabled
      when present) → save logits + labels → compute accuracy, F1, confusion
      matrix → update model card.
      \item Arrows: linear flow, annotate produced artifacts.
      \item Explanation: how evaluation outputs are derived and stored.
    \end{itemize}
\end{enumerate}

\vspace{1em}
\noindent All implementation references above (file paths and symbol names)
are exact and present in the repository. This chapter preserves the factual
implementation details whilst providing expanded methodological justification
and a formal academic presentation suitable for a thesis.

\end{document}

